% Hand-written discussion for make_report.py. It interprets the numbers of one run; make_report.py flags it
% when the run it is built with is a different one.
% written for run: 2026-10-13T09:39:49
\section{Discussion}
Reprojection numbers below are for the left (reference) camera; stereo numbers use both cameras.

\subsection{T1: the in-air calibration cannot be used underwater}
Used unchanged, the in-air calibration gives a held-out error of 8.6~px in the left camera, 30 times that
of the underwater fits. It triangulates the 250--450~mm scale bars 29.5~mm (about 8\%) short and the 145~mm
target tags 6--14~mm short, the error growing with distance (Table~\ref{tab:target}). Both follow from the
focal length. Water multiplies it by about $n$: the underwater fits give $f_x \approx 953$~px against
719~px in air, a ratio of 1.33. In a stereo pair the depth is $Z = fB/d$, so a focal length that is 25\%
too short puts every point 25\% too close. The sideways position $X = (u - c_x)B/d$ does not depend on $f$
at all. A tilted board or tag is therefore squashed along the viewing direction only, and its measured
lengths come out short by an amount that depends on its tilt. The quiver plot (Figure~\ref{fig:fig_quiver})
shows that the misfit is mostly radial, as a focal-length error is, and largest at the edges.

The held-out error of T1 is smaller far away (1.7~px at 1--1.5~m against 12.9~px near). This is not because
the model is better there. Far boards fill only the middle of the image, where a focal-length error can be
taken up by moving the board, because each board gets its own fitted pose. The triangulated lengths, which
use no board pose, are still about 26~mm short at 1--1.5~m.

\subsection{T2: Pinax fixes most of it, but not the edges}
Adding Snell's law at a flat port (Pinax) brings the held-out error down to 1.9~px and the scale-bar bias
from $-29.5$ to $+1.8$~mm, without fitting anything in water. That confirms refraction is what breaks T1.
Pinax is still about 7 times worse than the underwater fits, though, and its error grows steeply towards the
image edge: 9.5~px in the outer zone, against 0.47~px for Brown-Conrady
(Figure~\ref{fig:fig_radial}). The ray-angle comparison (Figure~\ref{fig:fig_pinax_mapping}) shows why.
Near the axis, Pinax and the in-water fit agree. Further out, Pinax's underwater rays fall short, by about
$0.5^\circ$ at 600~px from the centre and up to $2.7^\circ$ in the corners. A focal-length difference alone
would explain only $0.15^\circ$ of that. The Aquagon is a corrective water-contact optic, not a camera
behind a flat window, so the flat-port model over-corrects the off-axis rays it was designed to fix.

\subsection{T3 vs T4: the underwater fits are equivalent}
Brown-Conrady and Kannala-Brandt fitted in water give the same held-out error (0.29 and 0.28~px), and
every paired difference is within noise (Table~\ref{tab:paired}). Kannala-Brandt is slightly lower at the
image edge (0.43 against 0.47~px, also within noise).
Brown-Conrady triangulates slightly better (1.28 against 1.44~mm length RMS, within noise). On the separate
target both measure the tags to within 0.15~mm of each other at every distance. The correlation matrices
(Figure~\ref{fig:fig_correlation}) show that Kannala-Brandt's four distortion coefficients are almost
interchangeable ($|C|$ = 0.81--0.99): the data constrain its distortion curve, not the individual
coefficients. Brown-Conrady's $k_1, k_2, k_3$ behave the same way ($|C|$ up to 0.97). None of this harms
accuracy, but it means neither model's coefficients should be read physically. With the Aquagon's mild
distortion, the extra flexibility of the fisheye model buys nothing here. Brown-Conrady, the standard
pipeline model, is adequate.

\subsection{T5 and the close-range checks: the camera behaves as a single-viewpoint camera}
Brown-Conrady plus an axial ray shift is a non-single-viewpoint model that contains T3 as a special case. If
the Aquagon's rays did not meet in one point, T5 should predict held-out boards better, especially when
trained close and tested far. It does not. The held-out reprojection error is unchanged to 0.001~px, and
every paired difference with T3, and with T3 trained close-range only, is within noise. The scale-bar RMS is
0.17~mm lower with T5, also within noise. The fitted spread of where rays cross the axis is small, 0.5~mm at
$37^\circ$ off-axis, and varies between splits from 0.3 to 0.7~mm, which is what a weakly determined
parameter looks like, not a physical port offset.

Training Brown-Conrady on boards at 0.2--0.8~m only and scoring it out to 1.45~m (the \emph{close-range}
rows) costs nothing measurable: all of its differences from the full fit are within noise. A model that
transfers from near to far like this is what a single-viewpoint camera allows. A camera whose rays did not
meet in one point would show error growing with distance outside the training band. It does not
(Figure~\ref{fig:fig_distance}).

\subsection{The separate target: the remaining error is not the camera model}
On the WaterComparison tags, the five underwater-fitted models (T3, T4, T5 and the two close-range
variants) agree with each other to within about 0.3~mm at every distance. All of them measure the 145~mm
tags $-0.5$~mm short at 50~cm, $-1.5$~mm at 100~cm and $-3.5$ to $-3.8$~mm at 150--200~cm (about 2.5\%).
Because models with very different assumptions, including the non-single-viewpoint one, give the same error
on the same pixels, the error comes from the pixels, not from the camera model. The likely sources are:
\begin{itemize}
\item the corner localisation on small, hazy tags (at 2~m a tag is about 67~px wide, and 1~px of corner
bias is about 2\% of its size);
\item the stereo extrinsics, which are shared in form by all the models;
\item the target itself.
\end{itemize}
An earlier test with different corner refinements showed this error changes by about 2~mm at 1.5--2~m with
the refinement method alone. The 250~cm label has a single tag with inconsistent edges and should not be
used.

\subsection{What the covariance says, and what it does not}
The formal standard deviations are small: $\pm0.3$~px on $f_x$, $\pm0.1$~px on the principal point. The fit
noise they assume ($\hat\sigma = 0.19$~px for Brown-Conrady) is, however, lower than the held-out error
(0.29~px): pixel errors are not independent white noise, because corners of one tag and frames of one
sweep share their errors. The formal uncertainties are therefore optimistic by roughly 1.5 times, and
significance tests built on them over-claim. T5 is the clear example. Its shift coefficient
$a_2 = 0.95 \pm 0.14$~mm is 7 formal standard deviations from zero, yet across the 28 splits it ranges from
0.4 to 2.0~mm (standard deviation 0.39~mm, almost three times the formal 0.14~mm) and it does not improve
prediction at all. The held-out comparisons in this report, which use
recordings the fit never saw, are the test to rely on. The covariance is useful for something else: it shows
which parameters trade off ($f_x$/$f_y$ always, the radial coefficients among themselves, and $f$ with $k_1$
in the in-air fit).

\subsection{Conclusions}
\begin{enumerate}
\item An in-air calibration cannot be reused underwater: it measures lengths about 8\% short.
\item Physics-based flat-port compensation (Pinax) recovers most of the error without underwater data, but
not at the edges, because the Aquagon is not a flat port. It is about 7 times worse than calibrating in water.
\item Calibrating in water with the standard Brown-Conrady model is as good as any model tested
(0.28--0.29~px held-out, about 1.3~mm length RMS on 250--450~mm bars). Kannala-Brandt and the
non-single-viewpoint extension add nothing measurable.
\item Within the data's range (0.23--1.42~m, upper 71\% of the image), the camera behaves as a
single-viewpoint camera: a nested non-single-viewpoint model and near-to-far transfer both find nothing.
\item The 2--3\% tag-size error beyond the calibration range is common to every model, so it is not a
camera-model effect. Corner localisation on the far, hazy tags is the most likely cause.
\end{enumerate}

\subsection{Limits}
The boards never reached below $v \approx 850$~px or beyond 1.45~m, so the bottom of the image and longer
ranges are extrapolated by every model. The WaterComparison target has one stereo pair per distance and
4--6 tags each. The water type and whether the focus changed between the air and water sessions are not
known; both affect T1 and T2 only. The axial model is one particular non-single-viewpoint form (symmetric
about the axis). A general per-pixel ray model was not tested, because the image coverage and depth range
would leave it poorly constrained.
